\section{Step 11 Add before narrowing the result}\label{pcm-s11}

Use \texttt{addu} to form the exact sum in \texttt{t0}. Its word operation is sufficient here because the two signed halfwords sum to a value between -65536 and 65534, within signed thirty-two-bit range. The u suffix suppresses an overflow exception; it does not reinterpret the inputs as unsigned mathematical samples. Do not store or discard high sum bits yet.

\Needspace{12\baselineskip}
\subsection{Write and inspect}
Replace the temporary branch and nop in \texttt{exact\_sum}. Keep the surrounding labels and region markers.

\Needspace{4\baselineskip}
\MARSCodeLabel{MIPS assembly}
\begin{minted}[samepage=true]{mars-idea}
    addu $t0, $t0, $t1
\end{minted}
Use recipe \texttt{s11-1}, \texttt{s11-2} from the companion construction inputs. Pause at \texttt{mix\_ready} before comparing this partial state.

\subsection{Expected observation}
A=30000 and B=10000 produce exact sum 40000. A=-30000 and B=-10000 produce -40000.

\Needspace{11\baselineskip}
\subsection{Check your understanding}
Why does this word addition preserve the exact sum?

\begin{itemize}
  \item \textbf{A} Halfword inputs can never clip.
  \item \textbf{B} Their complete sum fits signed thirty-two bits.
  \item \textbf{C} \texttt{addu} automatically saturates at 32767.
\end{itemize}
